\documentclass[10pt,a4paper]{amsart}
\usepackage[margin=19mm]{geometry}
\usepackage{amsmath,amssymb,mathtools}
\usepackage[scr=rsfs,cal=euler]{mathalfa}
\usepackage{xfrac}
\usepackage[unicode,hidelinks,bookmarks=false]{hyperref}
\newcommand{\ie}{i.e.\ }
\newcommand{\cf}{cf.\ }
\newcommand{\resp}{respectively\ }
\newcommand{\Subsection}{\subsection}
\providecommand{\nobreakdash}{\nobreak\hbox{-}}
\setlength{\emergencystretch}{2em}
\multlinegap0pt
\usepackage{bm}
\usepackage[shortlabels]{enumitem}
\usepackage{float}
\usepackage{amssymb}
\usepackage{dsfont}
\usepackage{xcolor}
\usepackage{tikz}
\newtheorem{lemma}{Lemma}
\newtheorem{definition}{Definition}
\newtheorem{corollary}{Corollary}
\newtheorem{theorem}{Theorem}
\newcommand{\lam}{\lambda}
\title{Core abaci and Diophantine equations I: fundamental weight}
\author{Yanbo Li, Jiansheng Zhang, Shasha Zhu}
\date{Source: arXiv:2604.23311v2 (CC BY 4.0)}
\begin{document}
\maketitle

\noindent\emph{Source and licence.} Core abaci and Diophantine equations I: fundamental weight. Version arXiv:2604.23311v2 (CC BY 4.0), distributed by arXiv under the Creative Commons Attribution 4.0 International licence, http://creativecommons.org/licenses/by/4.0/, as recorded in the arXiv licence metadata for this exact version and shown on the abstract page https://arxiv.org/abs/2604.23311v2. Full licence notice: \texttt{licenses/arxiv-2604.23311-CC-BY-4.0.txt}.\par\smallskip

\noindent\textit{Editorial scope.} Counted unit: the article's theorem theorem (source0.tex lines 3152--3161). The article's complete original proof of it is delivered with it (source0.tex lines 3163--3243, one \texttt{proof} environment of 81 lines). No mathematics is written, repaired or re-derived: the statement and the proof are the article's own lines, in the article's own notation, and no retained line is substituted or re-flowed.  Retained uncounted as the source-local context the proof needs: lines 2700--2713; lines 2767--2771; lines 2827--2831.  Nothing was omitted from any retained span, so the package carries no omission record.  No retained line contains a citation, so the wrapper carries no bibliography and no bibliography entry.  No retained line contains a figure environment, an image-inclusion command or drawing code, so no image is delivered and the package declares no asset dependency.  Editorial formatting only, in addition: an AMS \texttt{amsart} preamble that restates the article's own declarations and macros used by the retained lines, namely the environments \texttt{lemma}, \texttt{definition}, \texttt{corollary}, \texttt{theorem} on separate counters, together with the standard packages those lines need.  The retained environments print in this document as Lemma 1, Definition 1, Corollary 1, Theorem 1 in order of appearance, whereas the article prints the same items with the numbering of its own full document.  The complete native source of the article as distributed in the arXiv e-print for this exact version is bundled unchanged as \texttt{sources/199154/source0.tex}.
\begin{lemma}\label{6.1}
A solution $t$ of an affine type Diophantine equation $x_1^2+x_2^2+\dots+x_l^2=\mathtt{a}N+\mathtt{b}$
is parameterized by a core abacus $(\lam, j)$ if and only if for $i=1, \dots, l$
\begin{enumerate}
\item [{\rm (1)}] $\frac{t_i+\mathrm{c}_i}{\mathtt{k}}\in\mathbb{Z}$
and the number of odd $\frac{t_i+\mathrm{c}_i}{\mathtt{k}}$ is $j$, for type
$C_l^{(1)}$ with charge $j$, 
or type $A_{2l-1}^{(2)}$, $A_{2l}^{(2)}$, $B_l^{(1)}$, $D_l^{(1)}$ or $D_{l+1}^{(2)}$ with charge $j$ satisfying $\frac{a_j^\vee}{a_0^\vee}=2$;
\item [{\rm (2)}] $\frac{t_i+\mathrm{c}_i}{\mathtt{k}}\in\mathbb{Z}$, for type
  $A_{2l-1}^{(2)}$, $B_l^{(1)}$ or $D_l^{(1)}$ with charge $j$ being $0$ or $1$, or type $A_{2l}^{(2)}$ or $D_{l+1}^{(2)}$ with charge $j=0$;
\item [{\rm (3)}] $\frac{t_i+\mathrm{c}_i}{\mathtt{k}}\in\mathbb{Z}+\frac{1}{2}$, for type
$D_l^{(1)}$ with charge $l-1$ or $l$, or type $B_{l}^{(1)}$ or $D_{l+1}^{(2)}$ with charge $l$.
\end{enumerate}
\end{lemma}

\begin{definition}\label{6.3}
Let $k\in \mathbb{Z}$ and $(b_1, \dots, b_l)\equiv (r_1, \dots, r_l)$, $(n_1, \dots, n_l)\equiv (r_1', \dots, r_l') \pmod k$, 
where $b_i, n_i\in \mathbb{Z}$ for $i=1, \dots, l$. Define $(b_1, \dots, b_l) \sim_k (n_1, \dots, n_l)$ if there exists 
$\sigma\in \mathfrak{S}_l$ such that $r_i=\sigma(r_i')$ or $r_i+\sigma(r_i')=k$ for each $i=1, \dots, l$.
\end{definition}

\begin{corollary}\label{6.5}
Let ${\rm O}_v$ and ${\rm O}_w$ be two solution orbits of an affine type $X$ 
Diophantine equation $x_1^2+x_2^2+\dots+x_l^2=\mathtt{a}N+\mathtt{b}$ with $[v]=[w]$.
Then the number of solutions being parameterized by core abaci in ${\rm O}_v$ is equal to that in ${\rm O}_w$.
\end{corollary}

\begin{theorem}
Given a nonnegative integer $N$, we have
$$|\mathcal{C}_j^X(N)|=
\begin{cases}
  \frac{1}{96}r_4(8N+6)=\frac{1}{4}\sum\limits_{d\mid 8N+6, 2\nmid d}d,\, & \mbox{\rm if } X=B_4^{(1)}, j=2, \ N \,\ \text{\rm is even}; \\
  \frac{1}{192}r_4(8N+6)=\frac{1}{8}\sum\limits_{d\mid 8N+6, 2\nmid d}d,\,\, & \mbox{\rm if } X=B_4^{(1)}, j=2, \  N \,\ \text{\rm is odd}; \\
  \frac{1}{24}r_4(6N+2)=\sum\limits_{d\mid 6N+2, 2\nmid d}d, & \mbox{\rm if } X=D_4^{(1)}, j=2,  \  N \,\ \text{\rm is odd}.
\end{cases}
$$
\end{theorem}

\begin{proof}
Let us consider affine type $B_4^{(1)}$ first. Note that the equations 
of affine $B_4^{(1)}$ type with charge $2$ are of the form $x_1^2+x_2^2+x_3^2+x_4^2=8N+6$.
For a given nonnegative integer $N$, let $t=(t_1, t_2, t_3, t_4)$ be a solution of the equation.
It is easy to check that two of $t_1, t_2, t_3, t_4$ are odd and ${\rm O}_{t}$ is in a unique $\sim_8$-equivalence class
$[(0, 1, 1, 2)]$, $[(0, 1, 2, 3)]$, $[(0, 2, 3, 3)]$, $[(1, 1, 2, 4)]$, $[(1, 2, 3, 4)]$ or $[(2, 3, 3, 4)]$ in the sense of Definition \ref{6.3}.
By Corollary \ref{6.5}, we only need to take one solution orbit in each equivalence class to compute $|{\rm O}_t^{\rm p}|$.
If $t\in [(0, 1, 1, 2)]$, without loss of generality, assume that $t_1=8k_1$, $t_2=8k_2+1$, $t_3=8k_3+1$ and $t_4=8k_4+2$. 
Then according to the definition of a solution orbit and Lemma \ref{6.1}, one can get by routine calculation 
that $|{\rm O}_t|$ and the number of solutions in ${\rm O}_t$ parameterized by core abaci are as follows.
$$(|{\rm O}_t|, |{\rm O}_t^{\rm p}|)=
\begin{cases}
		(96,\,1),  & {\rm if \,\,}  \text{ $k_1=0$, $|t_1|=|t_2|$}; \\
		(192,\,2),  & {\rm if \,\,}  \text{ $k_1\neq 0$, $|t_1|=|t_2|$};\\
		(192,\,2),  & {\rm if \,\,}  \text{ $k_1= 0$, $|t_1|\neq|t_2|$};\\
		(384,\,4),  & {\rm if \,\,}  \text{ $k_1\neq0$, $|t_1|\neq|t_2|$}.
\end{cases}$$
Similarly, if $t\in [(0, 1, 2, 3)]$, the corresponding datum is
$$(|{\rm O}_t|, |{\rm O}_t^{\rm p}|)=
\begin{cases}
	(192,\,1),  &  {\rm if \,\,} \text{ $k_1=0$}; \\
	(384,\,2),  &  {\rm if \,\,} \text{ $k_1\neq 0$.}
\end{cases}$$
If $t\in [(0, 2, 3, 3)]$, we have
$$(|{\rm O}_t|, |{\rm O}_t^{\rm p}|)=
\begin{cases}
		(96,\,1),  & {\rm if \,\,}  \text{ $k_1=0$, $|t_3|=|t_4|$}; \\
		(192,\,2),  & {\rm if \,\,}  \text{ $k_1\neq 0$, $|t_3|=|t_4|$};\\
		(192,\,2),  & {\rm if \,\,}  \text{ $k_1= 0$, $|t_3|\neq|t_4|$};\\
		(384,\,4),  & {\rm if \,\,}  \text{ $k_1\neq0$, $|t_3|\neq|t_4|$}.
\end{cases}$$
If $t\in [(1, 1, 2, 4)]$, then
$$(|{\rm O}_t|, |{\rm O}_t^{\rm p}|)=
\begin{cases}
		(192,\,2),  & {\rm if \,\,}  \text{ $|t_1|=|t_2|$};\\
		(384,\,4),  & {\rm if \,\,}  \text{ $|t_1|\neq|t_2|$}.
\end{cases}$$
If $t\in [(2, 3, 3, 4)]$, then
$$(|{\rm O}_t|, |{\rm O}_t^{\rm p}|)=
\begin{cases}
		(192,\,2),  & {\rm if \,\,}  \text{ $|t_2|=|t_3|$};\\
		(384,\,4),  & {\rm if \,\,}  \text{ $|t_2|\neq|t_3|$}.
\end{cases}$$
Finally, if $t\in [(1, 2, 3, 4)]$, then $(|{\rm O}_t|, |{\rm O}_t^{\rm p}|)=(384, \,2)$.

If $N$ is even, then $8N+6\equiv 6 \pmod {16}$. This forces the squares of two odd numbers in $t$ simultaneously satisfying 
$x^2\equiv 1 \pmod{16}$, or $x^2\equiv 9 \pmod{16}$ by analysis of modulo 16. 
As a result, $t\in [(0, 1, 1, 2)]$, $[(0, 2, 3, 3)]$, $[(1, 1, 2, 4)]$ or $[(2, 3, 3, 4)]$.
This implies that $|{\rm O}_t|/ |{\rm O}_t^{\rm p}|=96$ and thus $|\mathcal{C}_j^X(N)|=\frac{1}{96}r_4(8N+6)$.
On the other hand, for $N$ being odd, the squares of two odd numbers in $t$ 
satisfy that one is of the form $16n_1+1$ and the other one $16n_2+9$.
That is, $t\in [(0, 1, 2, 3)]$ or $[(1, 2, 3, 4)]$ and hence $|{\rm O}_t|/ |{\rm O}_t^{\rm p}|=192$. 
The proof of affine type $B_4^{(1)}$ is complete.

We now consider affine type $D_4^{(1)}$. Note that the corresponding equations are of the form $x_1^2+x_2^2+x_3^2+x_4^2=6N+2$. 
If $N$ is odd, then $4|6N+2$. Based on this fact, we can obtain by direct calculation 
that a solution $t$ of the equation is in a $\sim_6$ equivalence class $[(0, 0, 2, 2)]$ or $[(1, 1, 3, 3)]$.
If $t\in [(0, 0, 2, 2)]$, without loss of generality, assume that $t_1=6k_1$, $t_2=6k_2$, $t_3=6k_3+2$ and $t_4=6k_4+2$. Then
$$(|{\rm O}_t|, |{\rm O}_t^{\rm p}|)=
\begin{cases}
	(24,\,1),  &  {\rm if }\,\, \text{ $k_1=k_2=0$, $|t_3|=|t_4|$}; \\
	(48,\,2),  & {\rm if }\,\, \text{ $k_1=k_2=0$, $|t_3|\neq|t_4|$};\\
	(96,\,4),  & {\rm if }\,\, \text{ $k_1k_2=0$, $k_1\neq k_2$, $|t_3|=|t_4|$};\\
    (192,\,8),  & {\rm if }\,\, \text{ $k_1k_2=0$, $k_1\neq k_2$, $|t_3|\neq|t_4|$};\\
	(96,\,4),  & {\rm if }\,\, \text{ $k_1k_2\neq0$, $|t_1|=|t_2|$, $|t_3|=|t_4|$}; \\
	(192,\,8),  & {\rm if }\,\, \text{ $k_1k_2\neq0$, $|t_1|=|t_2|$, $|t_3|\neq|t_4|$};\\
	(192,\,8),  & {\rm if }\,\, \text{ $k_1k_2\neq0$, $|t_1|\neq|t_2|$, $|t_3|=|t_4|$};\\
	(384,\,16),  & {\rm if }\,\, \text{ $k_1k_2\neq0$, $|t_1|\neq|t_2|$, $|t_3|\neq|t_4|$.}
\end{cases}
$$
If $t\in [(1, 1, 3, 3)]$, then
$$(|{\rm O}_t|, |{\rm O}_t^{\rm p}|)=
\begin{cases}
	(96,\,4),  & {\rm if }\,\, \text{ $|t_1|=|t_2|$, $|t_3|=|t_4|$}; \\
	(192,\,8),  & {\rm if }\,\, \text{ $|t_1|=|t_2|$, $|t_3|\neq|t_4|$};\\
	(192,\,8),  & {\rm if }\,\, \text{ $|t_1|\neq|t_2|$, $|t_3|=|t_4|$};\\
	(384,\,16),  & {\rm if }\,\, \text{ $|t_1|\neq|t_2|$, $|t_3|\neq|t_4|$.}
\end{cases}
$$
That is, $|{\rm O}_t|/|{\rm O}_t^{\rm p}|=24$ for all cases and the proof is complete.
\end{proof}

\end{document}
